\subsection{Exponentials with Base Less than One} 
Previously, we saw what the graph of an exponential looks like if \makeblank[1in]. What if \makeblank[1in]?
\begin{example}
Sketch the graph of $f(x)=\left(\tfrac{1}{2}\right)^x$
\begin{lpic}{}{10}
	\draw (1,-2) -- (5,-2);
	\draw (3,-1) -- (3,-9);
	\regtext{2}{-2}{$x$};
	\regtext{4}{-2}{$y$};
	\foreach \x in {-3,...,3}
		\regtext{2}{-\x-6}{$\x$};
\end{lpic}
\end{example}
What do we see in this shape?
\begin{itemize}
\item The domain is \makeblank.
\item The range is \makeblank.
\item We can also describe the end behavior:
\begin{itemize}
\item When $x\rightarrow+\infty$, $f(x)\rightarrow\makeblanksmall$.
\item When $x\rightarrow-\infty$, $f(x)\rightarrow\makeblanksmall$.
\end{itemize}
\item This graph also has a \makeblank[1in] asymptote at \makeblank[1in].
\end{itemize}
The key difference between $b>1$ and $b<1$ is the \makeblank.